% ==============================================================================
% sci_s3_motor_controller.tex — SCI/SICE tutorial Session 3:
% motor control and controller input implementation
% SCI/SICE チュートリアル セッション3: モータ制御とコントローラ入力の実装
% ==============================================================================

% ------------------------------------------------------------------
\begin{frame}{地図: 今ここ}
  \begin{columns}[c]
    \begin{column}{0.5\textwidth}
      \centering
      \resizebox{\columnwidth}{!}{\scisessionmap{3}}
    \end{column}
    \begin{column}{0.47\textwidth}
      \small
      S2 でセンサの値を読めるようになった。ここでは読んだ値を使って機体を動かす側（モータ・コントローラ）を作る
    \end{column}
  \end{columns}
\end{frame}

% ------------------------------------------------------------------
\begin{frame}{要点3つ}
  \begin{enumerate}
    \item モータは PWM（パルス幅変調）の Duty 比で駆動する。4基の配置と回転方向はトルクバランスで決まっている
    \item コントローラのスティック値は ESP-NOW で 14 バイトの \code{ControlPacket} として届く
    \item \api{ws::motor\_mixer(T,R,P,Y)} で手動ミキシングしても、フィードバックがない限り安定して飛べない。ここに次セッションの動機がある
  \end{enumerate}
\end{frame}

% ------------------------------------------------------------------
\begin{frame}{安全}
  \begin{alertblock}{モータを回す実習の注意}
    \begin{itemize}\small\setlength{\itemsep}{2pt}
      \item \warn{必ず机上}で行う。回転中はモータに手を近づけない
      \item \warn{Duty は 0.15 以下}に固定する。ファーム側では強制されないので数値を必ず確認する
      \item \warn{ARM} は機体ボタンの単クリックまたはコントローラで行う。コードの中で自動 ARM しない
      \item 異常な振動・音・回転が起きたら即座に \warn{DISARM}（機体ボタン再クリック）
    \end{itemize}
  \end{alertblock}

  \vspace{0.5em}

  {\small 止め方: 機体ボタンをもう一度クリック（DISARM）。バッテリー駆動での確認を基本とする}
\end{frame}

% ------------------------------------------------------------------
\begin{frame}{モータと PWM}
  \begin{columns}[T]
    \begin{column}{0.45\textwidth}
      \centering
      \resizebox{!}{0.62\textheight}{\input{../../_shared/tikz/pwm_waveform}}
    \end{column}
    \begin{column}{0.5\textwidth}
      \vspace{1em}
      \begin{itemize}\small
        \item LEDC（ESP32 のハードウェア PWM）でモータを駆動する
        \item Duty 0\% = 停止、100\% = 最大回転
        \item API: \api{ws::motor\_set\_duty(id, duty)}\\
              \small id=1--4, duty=0.0--1.0
      \end{itemize}
    \end{column}
  \end{columns}
\end{frame}

% ------------------------------------------------------------------
\begin{frame}{モータ配置}
  \begin{center}
    \resizebox{0.4\textwidth}{!}{\input{../../_shared/tikz/motor_layout}}
  \end{center}
  \begin{itemize}\small
    \item モータ ID: 1=FR（右前）、2=RR（右後）、3=RL（左後）、4=FL（左前）
    \item 対角のモータが同じ方向に回転する。M1・M3 = CCW、M2・M4 = CW（反トルクを打ち消す）
  \end{itemize}
\end{frame}

% ------------------------------------------------------------------
\begin{frame}[fragile]{実習 3: Duty をハードコードして回す}
  \vspace{0.15em}
  {\small 見るだけ ｜ \textcolor{sfblue}{\textbf{一緒に打つ}} ｜ 帰宅後に再現}
  \vspace{0.15em}

  \begin{columns}[T]
    \begin{column}{0.44\textwidth}
      \footnotesize
      \begin{enumerate}\setlength{\itemsep}{1pt}
        \item \code{sf lesson switch sci2026:3}
        \item 右のコードに書き換える
        \item \code{sf lesson build}
        \item \code{sf lesson flash}
        \item 機体ボタンを1回クリックして ARM
        \item 確認: FR モータだけが回る
      \end{enumerate}
    \end{column}
    \begin{column}{0.54\textwidth}
      \begin{sflisting}[basicstyle=\ttfamily\scriptsize]{user\_code.cpp}
void loop_400Hz(float dt) {
    ws::motor_set_duty(1, 0.10f);   // FR
    ws::motor_set_duty(2, 0.00f);   // RR
    ws::motor_set_duty(3, 0.00f);   // RL
    ws::motor_set_duty(4, 0.00f);   // FL
}
      \end{sflisting}
    \end{column}
  \end{columns}

  \vspace{0.15em}

  {\footnotesize \warn{注意}: 配布される模範解答（\texttt{--solution}）は起動時に自動 ARM する。必ず机上に置いた状態で起動すること（コードの中で自動 ARM しない）}
\end{frame}

% ------------------------------------------------------------------
\begin{frame}{コントローラ}
  \begin{center}
    \adjustbox{max width=0.92\textwidth, max height=0.78\textheight}{\input{../../_shared/tikz/controller_annotated}}
  \end{center}
\end{frame}

% ------------------------------------------------------------------
\begin{frame}{ESP-NOW と ControlPacket}
  \begin{center}
    \adjustbox{max width=0.92\textwidth, max height=0.68\textheight}{\input{../../_shared/tikz/espnow_dataflow}}
  \end{center}

  \vspace{0.3em}

  {\small コントローラとStampFlyは ESP-NOW で直接通信する。1パケット \code{ControlPacket} は \textbf{14バイト}（throttle/roll/pitch/yaw + flags + checksum）を約 \textbf{50\,Hz}（TDMA フレーム周期 20\,ms）で送信する}
\end{frame}

% ------------------------------------------------------------------
\begin{frame}{TDMA で 30 機を同時に飛ばす}
  \begin{block}{TDMA（時分割多重アクセス、Time Division Multiple Access）}
    \footnotesize 1つの無線チャンネルを短い時間枠（スロット）に分割し、機体ごとに送信タイミングをずらして電波の衝突を避ける方式
  \end{block}

  \vspace{0.2em}

  \begin{table}
    \centering\footnotesize
    \begin{tabular}{ll}
      \hline
      フレーム周期 20\,ms、スロット幅 2\,ms & $\to$ 1フレーム10スロット \\
      \textbf{1チャンネルあたり最大10機} & 非重複3チャンネル（1/6/11ch）で\textbf{理論上30機} \\
      \hline
    \end{tabular}
  \end{table}

  \vspace{0.1em}

  \begin{columns}[T]
    \begin{column}{0.48\textwidth}
      \begin{alertblock}{なぜ必要か}
        \scriptsize 教室で複数の送信機が無調整のまま同時に電波を出すと衝突し、機体が通信途絶と判断して自動着陸する
      \end{alertblock}
    \end{column}
    \begin{column}{0.5\textwidth}
      \begin{exampleblock}{スロットの割当}
        \scriptsize タッチメニューで \textbf{Device ID}（0--9）を手動設定し NVS（不揮発性メモリ）に保存（再起動後に反映）。ID=0 が20\,msごとにビーコンを送る「親機」、他はその受信時刻を基準に自分のスロットのみで送信する。ペアリングでは自動割当されない
      \end{exampleblock}
    \end{column}
  \end{columns}

  \vspace{0.1em}

  {\scriptsize 実績: DXH2026講座で10台を1/6/11chに4/3/3台で分散、Device IDを個別割当し同時運用した}
\end{frame}

% ------------------------------------------------------------------
\begin{frame}[fragile]{ペアリングとチャンネル}
  \begin{block}{ペアリング}
    \footnotesize コントローラの M5 ボタン（LCD パネル）を押しながら電源ON $\to$ StampFly のボタンを約2秒長押し $\to$ 両方のビープで完了。次回以降は自動接続
  \end{block}

  \vspace{0.3em}

  \begin{alertblock}{チャンネル設定}
    \footnotesize 教室内で複数機体が同時に飛ぶと混信する。\api{ws::set\_channel(ch)} を \code{setup()} で呼んで機体のチャンネルを 1/6/11 のいずれかにし、再ペアリングする。コントローラはペアリング時に機体のチャンネルへ自動追従する
  \end{alertblock}

  \vspace{0.3em}

  \begin{exampleblock}{自分で購入したコントローラを使う場合}
    \footnotesize ご購入直後のコントローラには出荷時点のファームウェアが入っているが、本チュートリアルで使う版とは異なる。最初に一度だけ、\code{sf flash controller} で書き込み直す必要がある
  \end{exampleblock}
\end{frame}

% ------------------------------------------------------------------
\begin{frame}{オープンループ制御}
  \begin{center}
    \resizebox{0.7\textwidth}{!}{\input{../../_shared/tikz/open_loop}}
  \end{center}

  \vspace{0.3em}

  \begin{table}
    \centering\small
    \begin{tabular}{ll}
      \hline
      \api{ws::rc\_throttle/roll/pitch/yaw()} & スティック値の取得 \\
      \api{ws::motor\_mixer(T, R, P, Y)}      & 4モータへの配分 \\
      \hline
    \end{tabular}
  \end{table}
\end{frame}

% ------------------------------------------------------------------
\begin{frame}{ミキサ行列}
  \begin{columns}[T]
    \begin{column}{0.45\textwidth}
      \centering
      \resizebox{0.95\columnwidth}{!}{\input{../../_shared/tikz/mixer_matrix}}
    \end{column}
    \begin{column}{0.52\textwidth}
      \vspace{0.3em}
      {\footnotesize
      Pitch+ $\to$ 前傾 $\to$ 前(M1,M4)強 + 後(M2,M3)弱\\[2pt]
      Roll+ $\to$ 右傾 $\to$ 右(M1,M2)弱 + 左(M3,M4)強}
      \vspace{0.3em}
      \begin{exampleblock}{意味}\footnotesize
        各式は「T/R/P/Yの指令を4モータのDutyにどう配分するか」を表す。式ごとの符号がモータ配置（motor\_layout）と対応する
      \end{exampleblock}
      \vspace{0.2em}
      {\footnotesize 実機ファーム（sf\_actuator）は $\tau/L$，$\kappa$ を使う物理量ベースのミキサ。演習の \api{ws::motor\_mixer} は共通スケール（$0.25/3.7$）の簡易版}
    \end{column}
  \end{columns}

  {\scriptsize $T$ は \api{ws::rc\_throttle()} の正規化スティック値で $[0,1]$，$R,P,Y$ は \api{ws::rc\_roll/pitch/yaw()} で $[-1,1]$。\api{ws::motor\_set\_duty} は Duty を $[0,1]$ にクランプする}
\end{frame}

% ------------------------------------------------------------------
\begin{frame}[fragile]{実習 4: スティックでモータを回す}
  \vspace{0.15em}
  {\small 見るだけ ｜ \textcolor{sfblue}{\textbf{一緒に打つ}} ｜ 帰宅後に再現}

  \begin{columns}[T]
    \begin{column}{0.44\textwidth}
      \footnotesize
      \begin{enumerate}\setlength{\itemsep}{0pt}
        \item \code{sf lesson switch sci2026:4}
        \item コントローラとペアリング
        \item \api{ws::motor\_mixer(T,R,P,Y)} を書く
        \item \code{sf lesson build}
        \item \code{sf lesson flash}
        \item 確認: スティックで回転差が出る
      \end{enumerate}
    \end{column}
    \begin{column}{0.54\textwidth}
      \begin{sflisting}[basicstyle=\ttfamily\scriptsize]{user\_code.cpp}
void loop_400Hz(float dt) {
    float t = ws::rc_throttle();
    float r = ws::rc_roll();
    float p = ws::rc_pitch();
    float y = ws::rc_yaw();
    ws::motor_mixer(t, r, p, y);
}
      \end{sflisting}
    \end{column}
  \end{columns}

  \begin{alertblock}{これでは飛べない理由}
    \scriptsize モータの個体差，機体のミスアライメント（推力軸のずれ・重心ずれ），電池電圧の低下といった外乱に対して何の補正もかからないため，スティック指令どおりの回転を保てない
  \end{alertblock}
\end{frame}

% ------------------------------------------------------------------
\begin{frame}{プラントの定義}
  \small 本セッションで作った経路（ミキサ $\to$ モータ $\to$ 機体）が、そのまま制御対象（プラント）の入出力になる

  \begin{columns}[T]
    \begin{column}{0.4\textwidth}
      \centering
      \vspace{0.5em}
      \resizebox{0.95\columnwidth}{!}{%
      \begin{tikzpicture}[
        block/.style={rectangle, draw=sfdark, line width=1.2pt, fill=sflight,
                      rounded corners=2pt, minimum width=22mm, minimum height=13mm,
                      font=\small},
        arr/.style={-{Stealth[length=2.5mm]}, line width=1pt, draw=sfdark},
      ]
        \node[block] (plant) {プラント};
        \coordinate (pin) at ($(plant.west)+(-12mm,0)$);
        \coordinate (pout) at ($(plant.east)+(12mm,0)$);
        \draw[arr] (pin) -- (plant.west)
          node[pos=0.3, above, font=\scriptsize] {Duty};
        \draw[arr] (plant.east) -- (pout)
          node[pos=0.7, above, font=\scriptsize] {角速度 $\omega$};
      \end{tikzpicture}}
    \end{column}
    \begin{column}{0.55\textwidth}
      \vspace{0.3em}
      \begin{itemize}\footnotesize\setlength{\itemsep}{3pt}
        \item \textbf{入力}: ミキサが出す4つの Duty（\api{ws::motor\_mixer(T,R,P,Y)} の R/P/Y トルク指令に相当）
        \item \textbf{出力}: \api{ws::gyro\_x/y/z()} が返す角速度の実測値
      \end{itemize}
    \end{column}
  \end{columns}

  \vspace{0.5em}

  \begin{exampleblock}{次セッションへ}
    \footnotesize S4 の PID 制御器がこの間にフィードバックループを作り，出力を目標値に近づける。伝達関数とゲイン $K_p$ の設計が S4「フィードバック制御の基礎」の主題である
  \end{exampleblock}
\end{frame}

% ------------------------------------------------------------------
\begin{frame}{理論とコードの対応表\later}
  \vspace{0.3em}

  {\small ここまで登場した理論用語・API・通信・実装トピックの対応を一覧できるようにまとめた早見表}

  \vspace{0.3em}

  \resizebox{\textwidth}{!}{%
  \begin{tabular}{llll}
    \hline
    \textbf{理論} & \textbf{コード} & \textbf{通信} & \textbf{トピック} \\
    \hline
    Duty比 $[0,1]$ & \api{ws::motor\_set\_duty(id,d)} & --- & \code{actuator\_motor} \\
    スティック指令 $T,R,P,Y$ & \api{ws::rc\_throttle()} 等 & \code{ControlPacket}（14B） & \code{command\_setpoint} \\
    ミキシング & \api{ws::motor\_mixer(T,R,P,Y)} & --- & --- \\
    \hline
  \end{tabular}}
\end{frame}

% ------------------------------------------------------------------
\begin{frame}{復習パス\later}
  \resizebox{\textwidth}{!}{%
  \begin{tabular}{lll}
    \hline
    \textbf{コマンド} & \textbf{実習} & \textbf{文書} \\
    \hline
    \code{sf lesson switch sci2026:3} & 実習 3 モータ制御 & \code{firmware/vehicle/docs/hardware\_init.md} \\
    \code{sf lesson switch sci2026:4} & 実習 4 コントローラ入力 & \code{protocol/spec/}（ControlPacket） \\
    \code{sf flash controller} & --- & \code{docs/commands/sf-flasher.md} \\
    \hline
  \end{tabular}}

  \vspace{0.5em}

  {\small すべて \code{sf lesson build} $\to$ \code{sf lesson flash} で実機に反映する}
\end{frame}

% ------------------------------------------------------------------
\begin{frame}{チェックポイント}
  \begin{alertblock}{確認事項}
    \begin{itemize}
      \item[$\square$] Duty とモータ配置・回転方向の対応を説明できる
      \item[$\square$] ControlPacket がどの経路（ESP-NOW）でどれだけの頻度で届くか説明できる
      \item[$\square$] オープンループ制御が「なぜ飛べないか」を自分の言葉で説明できる
    \end{itemize}
  \end{alertblock}

  \vspace{1em}

  \begin{block}{次のセッション}
    S4: フィードバック制御の基礎 --- PID による姿勢安定化（14:00 -- 15:00）
  \end{block}
\end{frame}
